\documentclass[12pt,a4paper]{report}

\usepackage[margin=1in]{geometry}
\usepackage{xcolor}
\usepackage{sectsty}
\usepackage{fancyhdr}
\usepackage{setspace}

\definecolor{titlemagenta}{RGB}{200,0,200}
\definecolor{titleorange}{RGB}{230,120,0}

% Editable student/course details.
\newcommand{\StudentName}{Shlok Patel}
\newcommand{\StudentID}{[Your ID]}
\newcommand{\CourseTitle}{Technical Writing}
\newcommand{\InstructorName}{[Course Instructor]}

% Chapter/section/subsection heading styles.
\chapterfont{\color{red}\bfseries\Huge}
\sectionfont{\color{titlemagenta}\itshape\Large}
\subsectionfont{\color{titleorange}\bfseries}

% Header from page 2 onwards, including chapter-opening pages.
\pagestyle{fancy}
\fancyhf{}
\fancyhead[L]{\StudentName}
\fancyhead[C]{\CourseTitle}
\fancyhead[R]{\today}
\renewcommand{\headrulewidth}{0.4pt}
\renewcommand{\footrulewidth}{0pt}
\fancypagestyle{plain}{%
  \fancyhf{}
  \fancyhead[L]{\StudentName}
  \fancyhead[C]{\CourseTitle}
  \fancyhead[R]{\today}
  \renewcommand{\headrulewidth}{0.4pt}
}

\begin{document}

% First page: title page without any header/footer.
\begin{titlepage}
  \thispagestyle{empty}
  \centering
  \vspace*{2cm}
  {\Huge\bfseries Introduction to \LaTeX\par}
  \vspace{2cm}
  {\Large \StudentName\par}
  \vspace{0.5cm}
  {\large ID: \StudentID\par}
  \vspace{0.5cm}
  {\large \CourseTitle\par}
  \vspace{0.5cm}
  {\large Instructor: \InstructorName\par}
  \vfill
  {\large \today\par}
\end{titlepage}

% All pages after the title page use the requested header.
\begin{abstract}
This document introduces LaTeX as a document preparation system.
It demonstrates document structure, mathematical typesetting,
tables, figures, citations, and presentation-related workflows.
\end{abstract}

\tableofcontents

\chapter{Introduction}
LaTeX separates document content from formatting instructions.
It is especially useful for technical reports and mathematical documents.

\section{Purpose}
This section introduces the purpose and structure of the report.

\subsection{Main Features}
LaTeX supports automatic numbering, cross-references and equations.

\chapter{Document Structure}
\section{Sections and Subsections}
The report class supports chapters, sections and subsections.

\subsection{Example}
A subsection appears below its parent section.

\chapter{Mathematics and Figures}
\section{A Sample Equation}
\[
E=mc^2
\]

\section{A Sample Figure}
Figures can be inserted using the graphicx package.

\chapter{Conclusion}
This report demonstrates the basic structure of a LaTeX report.

\end{document}
